\section{Fixed channels and the equality profiles}
\label{rig:section}

The pointwise relations of Proposition~\ref{comp:relations} still involve
an auxiliary operator $B$ that may depend on the point without a prescribed
measurable choice.  We now extract constant subspaces from those relations.
On each nonzero subspace, $B$ becomes a uniquely determined matrix and
satisfies an ordinary differential equation.  This will identify every
real equality potential; realification and a scalar calculation will then
complete the proof of Theorem~\ref{thm:main}.

\subsection{The channel subspaces are constant}

Assume first that $W$ is real and is a nonzero equality potential.  Use the
rows $a_i$, numbers $k_i$, and operator $\Lambda=\operatorname{diag}(k_i)$
from Proposition~\ref{comp:relations}.  For each distinct value $\alpha>0$
in the list $(k_i)$, let $n_\alpha$ be its finite multiplicity, and let
$G_\alpha$ be the $n_\alpha\times m$ matrix whose rows are the $a_i$ with
$k_i=\alpha$.  Each $G_\alpha$ has entries in $H^1(\R)$, and we use their
continuous representatives.

At every point in the common full-measure set of that proposition, choose
one of the operators $B$ it supplies.  Since $B$ commutes with $\Lambda^2$,
it preserves each of its eigenspaces.  Moreover,
\[
 AA^T=(\Lambda^2-B^2)^{1/r}
\]
commutes with both $B$ and $\Lambda^2$.  Writing $B_\alpha$ for the block
on the $\alpha^2$-eigenspace therefore gives
\begin{align}
 G_\alpha G_\beta^T&=0 &&(\alpha\ne\beta),
 \label{rig:orthogonal-blocks}\\
 G_\alpha'&=B_\alpha G_\alpha,
 & [B_\alpha,G_\alpha G_\alpha^T]&=0,
 \label{rig:block-first-order}\\
 \alpha^2I-B_\alpha^2&=(G_\alpha G_\alpha^T)^r,
 &\|B_\alpha\|_{\op}&\le\kappa.
 \label{rig:block-square}
\end{align}
These statements hold almost everywhere and require no compatibility
between the choices of $B$ at different points.

\begin{lemma}\label{rig:fixed-right-spaces}
For every $\alpha$, the subspace
\[
 \mathcal R_\alpha=(\ker G_\alpha(x))^\perp\subset\R^m
\]
is independent of $x\in\R$.  The nonzero spaces $\mathcal R_\alpha$ are
mutually orthogonal, and
\begin{equation}\label{rig:dimension-bound}
 \sum_\alpha\dim\mathcal R_\alpha\le m.
\end{equation}
In particular, only finitely many $G_\alpha$ are nonzero.  The potential
$W$ has a continuous representative, equal to zero on the orthogonal
complement of these spaces and equal to
$q(G_\alpha^TG_\alpha)^r$ on $\mathcal R_\alpha$.  With this
representative, every $G_\alpha$ is $C^2$.
\end{lemma}

\begin{proof}
Fix $v\in\R^m$.  Equation~\eqref{rig:block-first-order} implies
\[
 \|(G_\alpha v)'(x)\|\le\kappa\|G_\alpha(x)v\|
 \quad\text{for almost every }x.
\]
The function $G_\alpha v$ is locally absolutely continuous.  If it
vanishes at $x_0$, integration and Gronwall's inequality, first to the
right and then to the left of $x_0$, show that it vanishes everywhere.
Consequently $\ker G_\alpha(x)$ is independent of $x$, including points
outside the original full-measure set.  This proves the first assertion.

The products in \eqref{rig:orthogonal-blocks} are continuous, so their
almost-everywhere vanishing implies their vanishing everywhere.  Thus the
constant row spaces $\mathcal R_\alpha$ are mutually orthogonal.  Any
finite collection has total dimension at most $m$, proving
\eqref{rig:dimension-bound} and finiteness of the nonzero blocks.

On the full-measure set from Proposition~\ref{comp:relations},
$A^TA=\sum_\alpha G_\alpha^TG_\alpha$.  Each summand acts on its own
orthogonal space $\mathcal R_\alpha$, and only finitely many summands
are nonzero.  Equation~\eqref{comp:potential-limit} therefore yields
\begin{equation}\label{rig:continuous-potential}
 W=q\sum_\alpha(G_\alpha^TG_\alpha)^r
 \quad\text{almost everywhere}.
\end{equation}
The right-hand side is continuous and has the asserted restrictions;
we henceforth use it as the representative of $W$.  The row eigenfunction
equation \eqref{comp:row-eigenfunction} now reads
\begin{equation}\label{rig:block-second-order}
 G_\alpha''=\alpha^2G_\alpha-G_\alpha W
 \quad\text{in distributions}.
\end{equation}
Its right-hand side is continuous.  Integrating this identity twice shows
that the continuous representative of $G_\alpha$ is $C^2$.
\end{proof}

The row spaces already split the physical space $\R^m$ into finitely
many constant channels.  To determine the potential within one of them,
we also need a constant range for $G_\alpha$ in the eigenfunction index
space $\R^{n_\alpha}$.

\begin{lemma}\label{rig:fixed-left-spaces}
For each $\alpha$ with $\ell:=\dim\mathcal R_\alpha>0$, the range of
$G_\alpha(x)$ in $\R^{n_\alpha}$ is independent of $x$.  In constant
orthonormal bases of this range and of $\mathcal R_\alpha$, the
restriction of $G_\alpha$ is an invertible $C^2$ matrix
$D:\R\to\R^{\ell\times\ell}$.  The matrix
$B_0=D'D^{-1}$ is $C^1$ and satisfies, everywhere on $\R$,
\begin{equation}\label{rig:active-identities}
 B_0=B_0^T,\qquad
 \alpha^2I-B_0^2=(DD^T)^r>0,\qquad
 B_0'=r(B_0^2-\alpha^2I).
\end{equation}
\end{lemma}

\begin{proof}
Choose a constant orthonormal basis of $\mathcal R_\alpha$, and let
$E(x)$ be the matrix of the restriction of $G_\alpha(x)$ to this space.
It has full column rank $\ell$ at every point, since the kernel of
$G_\alpha$ is constant.  The projection onto its range is
\[
 P(x)=E(x)(E(x)^TE(x))^{-1}E(x)^T.
\]
The commutator in \eqref{rig:block-first-order} shows that $B_\alpha$
preserves the range of $G_\alpha G_\alpha^T$, which is precisely the
range of $E$.  Hence $E'=B_\alpha E$ has its columns in this range
almost everywhere.  Differentiating the projection gives
\begin{align*}
 P'={}&(I-P)E'(E^TE)^{-1}E^T\\
     &+E(E^TE)^{-1}(E')^T(I-P)=0
 \quad\text{almost everywhere}.
\end{align*}
Since $P$ is $C^1$, it is constant.

Choose a constant orthonormal basis of this range as well.  The resulting
square matrix $D$ is invertible everywhere and $C^2$.  Almost everywhere,
$D'D^{-1}$ is the restriction of $B_\alpha$ to its invariant range in
these coordinates.  Equations~\eqref{rig:block-first-order} and
\eqref{rig:block-square} give the symmetry and square identity in
\eqref{rig:active-identities} almost everywhere; continuity extends
them to every point.  Positivity is strict because $D$ is invertible.

On $\mathcal R_\alpha$, Equation~\eqref{rig:continuous-potential}
identifies the potential with $q(D^TD)^r$.  Equation~
\eqref{rig:block-second-order} consequently becomes
\[
 D''=\alpha^2D-qD(D^TD)^r
     =\alpha^2D-q(DD^T)^rD.
\]
The second equality follows, for example, by a singular-value
decomposition.  Differentiating $B_0=D'D^{-1}$ and using $q=r+1$
now gives
\[
 B_0'=D''D^{-1}-B_0^2
      =\alpha^2I-q(\alpha^2I-B_0^2)-B_0^2
      =r(B_0^2-\alpha^2I),
\]
as claimed.
\end{proof}

\subsection{The matrix Riccati equation separates the profiles}

Fix a nonzero block and the matrices of Lemma~\ref{rig:fixed-left-spaces}.
Diagonalize the symmetric matrix $B_0(x_*)$ at one point by a constant
orthogonal change of range coordinates.  The right-hand side of its
Riccati equation is a polynomial in $B_0$.  Uniqueness for this matrix
ordinary differential equation therefore preserves the diagonal form.
More explicitly, the strict inequality in \eqref{rig:active-identities}
places each initial diagonal entry in $(-\alpha,\alpha)$.  There is a
unique $y_j\in\R$ such that the scalar solution with that initial value is
\begin{equation}\label{rig:tanh}
 b_j(x)=-\alpha\tanh\bigl(r\alpha(x-y_j)\bigr),
 \qquad 1\le j\le\ell.
\end{equation}
The diagonal matrix with these entries solves the same initial-value
problem as $B_0$ on all of $\R$, so uniqueness gives equality everywhere.

Equation~\eqref{rig:active-identities} implies that $DD^T$ is diagonal
in these coordinates.  Since $D'=\operatorname{diag}(b_j)D$, each row
of $D$ is a positive scalar multiple of its value at $x_*$.  None of
these rows is zero, and their directions are mutually orthogonal because
$DD^T$ is diagonal.  Taking their normalized directions as a constant
orthonormal basis of $\mathcal R_\alpha$ makes $D$ diagonal with positive
entries.  The potential in this fixed basis is therefore diagonal, with
entries
\begin{equation}\label{rig:profiles}
 q\bigl(\alpha^2-b_j(x)^2\bigr)
 =(r+1)\alpha^2\operatorname{sech}^2\bigl(r\alpha(x-y_j)\bigr).
\end{equation}
Combining these bases over the finitely many nonzero spaces
$\mathcal R_\alpha$, and adding a basis of their orthogonal complement,
proves the required necessity for real $W$.  There are at most $m$
nonzero profiles by \eqref{rig:dimension-bound}.  The zero potential
has the same description with no nonzero profile.

\subsection{Complex Hermitian potentials}

For a complex Hermitian potential $W$, take the realification
$\widetilde W$ from \eqref{bound:realification}.  The proof of
Proposition~\ref{bound:sharp} shows that it is equivalent to
$W\oplus\overline W$ by a constant unitary matrix and that both the trace-power
integral and the negative spectral moment are doubled.  Thus equality
for $W$ implies equality for the real potential $\widetilde W$.

The real result provides a constant unitary matrix $V$ such that
\[
 W(x)\oplus\overline W(x)
   =V\operatorname{diag}(f_1(x),\ldots,f_{2m}(x))V^*
 \quad\text{almost everywhere},
\]
where each $f_j$ is either a profile of the form \eqref{rig:profiles} or
zero.  Let $P$ be the projection onto the first summand $\C^m$ and put
$\widehat P=V^*PV$.  Since $P$ commutes with $W\oplus\overline W$,
\[
 (f_i(x)-f_j(x))\widehat P_{ij}=0
 \quad\text{almost everywhere}.
\]
Consequently $\widehat P_{ij}=0$ whenever the functions $f_i,f_j$ are
not equal almost everywhere.  Partition the indices into groups of
identical functions, including the possible zero function.  The range
of $\widehat P$ splits as the orthogonal sum of its intersections with
these coordinate groups.  On each group the diagonal potential is a
scalar multiple of the identity, so any constant orthonormal basis of
that intersection diagonalizes its restriction.  Transporting these
bases by $V$ to the first summand gives a constant unitary basis of
$\C^m$ in which $W$ has the asserted profiles.  This proves necessity
for complex Hermitian potentials.

\subsection{Each listed potential attains equality}

It remains to verify attainment, including the exact spectral moment,
without appealing to a classification of scalar optimizers.  Fix $a>0$
and $x_0\in\R$, and set
\[
 V(x)=(r+1)a^2\operatorname{sech}^2\bigl(ra(x-x_0)\bigr),
 \qquad H=-\partial_x^2-V.
\]
The first-order factorization used in the commutation argument
of~\cite[Section~II]{BenguriaLoss2000} also explains why this profile has
only one negative eigenvalue.  Put
\[
 w(x)=a\tanh\bigl(ra(x-x_0)\bigr),
 \qquad Q=\partial_x+w,\qquad \operatorname{Dom}Q=H^1(\R).
\]
Then $Q^*=-\partial_x+w$ on $H^1(\R)$.  Since $w$ and $w'$ are bounded,
both product domains are $H^2(\R)$, and direct differentiation yields
\begin{align}
 Q^*Q&=H+a^2,
 \label{rig:factorization}\\
 QQ^*&=-\partial_x^2+a^2
       +(r-1)a^2\operatorname{sech}^2\bigl(ra(x-x_0)\bigr)
       \ge a^2.
 \label{rig:partner}
\end{align}
Here the last inequality uses $r>1$.  The equation $Q\psi=0$ has the
one-dimensional $L^2$ solution space
\[
 \ker Q=
 \operatorname{span}\left\{
    \operatorname{sech}^{1/r}\bigl(ra(x-x_0)\bigr)
 \right\}.
\]
Thus $-a^2$ is a simple eigenvalue of $H$, and
\eqref{rig:factorization} excludes eigenvalues below it.  If
$H\phi=\lambda\phi$ with $-a^2<\lambda<0$, put
$\mu=\lambda+a^2$ and $g=Q\phi$.  Then
$\|g\|_2^2=\mu\|\phi\|_2^2>0$.  Moreover,
$Q^*Q\phi=\mu\phi\in\operatorname{Dom}Q$, so
$g\in\operatorname{Dom}(QQ^*)$ and $QQ^*g=\mu g$.  This contradicts
\eqref{rig:partner}, since $\mu<a^2$.  By Lemma~\ref{pre:spectral},
all negative spectrum is discrete.  Hence $-a^2$ is the sole negative
eigenvalue.

Finally, $2p-1=2\gamma$ and the substitution $s=\tanh y$ give
\begin{align*}
 \int_\R V(x)^p\,dx
 &=\frac{(r+1)^p}{r}a^{2\gamma}
     \int_\R\operatorname{sech}^{2p}y\,dy\\
 &=\frac{(r+1)^p}{r}a^{2\gamma}
     \int_{-1}^{1}(1-s^2)^{p-1}\,ds
  =\frac{b_\sigma}{D_\sigma}a^{2\gamma},
\end{align*}
where $(r+1)^p/r=D_\sigma^{-1}$ and $p-1=\sigma$.  Consequently
\[
 \operatorname{Tr}(H)_-^\gamma
   =a^{2\gamma}
   =C_\gamma\int_\R V^p.
\]
The zero potential contributes zero to both sides.  Finite orthogonal
direct sums add the spectral moments and trace-power integrals, and a
constant unitary conjugation leaves both unchanged.  Every potential
listed in Theorem~\ref{thm:main} therefore attains equality, completing
its proof.
